\section{Related Work}
\subsection{Long Range Propagation}

\paragraph{Attention Based Method}
Transformers \cite{vaswaniAttentionAllYou2023} model the depenencies between
feature sequence using self-attention mechansim. Graph nodes can be represented
as a seqeunce of tokens, such that the computed attention matrix is treated
like a new weighted adjacency matrix. The problem with this approach is that
the transformer is not aware of the graph structure. Works like Graphormer
\cite{yingTransformersReallyPerform2021a} use positional and structural
embeddings to inject the graph structure. GraphGPS
\cite{rampasekRecipeGeneralPowerful2023} computes global attention to learn
depenencies between distant nodes, and local MPNN layer to learn depenencies
between local nodes.

\paragraph{Implicit GNN}
\cite{bakerImplicitGraphNeural,bakerMonotoneOperatorTheoryInspired2024}
IGNN achive long-range propgation by solving the equalibrium equation $X =
    f_\theta(X,A)$, where $A$ is the   adjacency matrix and $X$ are the node
signals. For intance, the original IGNN model \cite{guImplicitGraphNeural2021}
finds $X$ such that $X = \sigma(WXA + b)$. They solve it by repeatadly applying
the equation using the same learned weight matrices until convergance, allowing
information to propagate through long distances without the need for deep
propagation.

\subsection{Learning Preconditioners}
\cite{liLearningPreconditionerConjugate2023} learnes a preconditioner $P =
f_\theta(A,b)$ for the discritzed PDE linear system $Ax = b$. They treat $A$ as
an adjacency matrix, having $A_{ij}$ as an edge feature and $b_{i}$ as a node
feature. They encode the graph into a hidden represnentation using MLP encoder,
then they follow it by multiple MPNN layers and pass into a decoder that
predicts a matrix $M$ that is used to construct the preconditioner $P$. In
order to gaurantee symmetry they learn $LL^T$.

\cite{boothNeuralAccelerationIncomplete} contructs a two-layer sparse linear
neural network whose weights are the entries of the incomplete Cholesky factor
$L$. The sparisty pattern is fixed in advance. The first layer applies $L^\top$
while the second applies $L$.

\cite{hausnerNeuralIncompleteFactorization2024} learns a lower-triagonal matrix
$L_\theta(A)$ with strictly positive diagonal, where the final preconditioner
becomes $L_\theta(A)L^\top_\theta(A)$. They do that by firstly propagating
information using MPNN layers of the lower-triagonal part of A, then they do
the same with the upper triagonal part of A. Finally the edge embeddings become
the entries of a sparse matrix triagonal factor.

% \cite{hausnerLearningIncompleteFactorization} learns a preconditioner for a
% non-singluar matrix $A$. Their approach differs from previous approcahes in
% that they do not require SPD or symmetic matrix. They learn a preconditioner of
% the form $P = LU$. They do that by using multiple message passing layers, with an
% additional edge feature to tell if the edge should belong to the upper or
% lower triangle matrix.
%
\cite{chenGraphNeuralPreconditioners2025} uses a GCN based method for learning
a network $M$ that acts as an approximation to the inverse. They do that by
first scaling the input vector $b$ and pass it into an encoder for a hidden
represnentation, then they pass it into multiple GCN layers, decode it back to
its original dimenstion and rescale it.

Most methods we've looked at directly learn a preconditioner.
\cite{trifonovLearningLinearAlgebra2025} proposed to use GNNs to learn a
correction for the given matrix. Given a lower triagonal IC factor $L$ of $A$,
their goal is to learn $L(\theta) = L + \alpha\ \mathrm{GNN} (\theta,L)$. They
use simple encodig, propagation using MPNN and decoding back.

% \cite{sapplConvolutionalNeuralNetworkdriven2026} used a CNN architcture in
% order to learn a lower triagonal matrix $L$ for a preconditioner $P = LL^T$.
%
% Algebric multigrids (AMG) are also used for learning preconditioners. AMG
% contructs a hierarchy of coarse represnentations directly from the given
% matrix. \cite{luzLearningAlgebraicMultigrid2020} uses GNN based architcture to
% learn the prolongation weights between the coarse and fine levels.
% \cite{liNeuralPreconditioningOperator2025} combined algebric multigrid
% principles with a transformer based architicture to learn a multi-scale
% preconditioning operator.
%
% \cite{brardaFactoredSparseApproximate2026} learns a fixed pattern sparse
% approximation inverse  $M = GG^\top$. The differnece with from other methods is
% that they minimize spectral objectives of $M$ rather than Frobenius norm.
%
%
% Generative methods where used by \cite{liGenerativeModelingSparse2024}. Given
% an SPD matrix $A$, they encode it using a GAT based encoder, while also passing
% the exact inverse $A^{-1}$ into a CNN encoder. The outputs are then
% concatanated and mapped into the latent space. A decoder then samples from the
% latent space and generates a sparse matrix $R$ that behaves like the inverse
% factor, meaning that $R^\top A R \approx I$.

